% =====================================================================================
% Methods -- consolidated. This file holds TWO top-level sections:
%   \section{User Study Methodology}                  (participants .. procedure)
%   \section{Analytical Strategy and Evaluation Metrics}  (telemetry .. modeling)
%
% It is the single source of truth for both. Earlier drafts pasted into chat or other
% files are stale; edit here.
%
% The methodology section reflects the following decisions, all of which reverse
% earlier drafts:
%   * n = 24, not 40. Sensitivity analysis replaces the a priori power calculation.
%   * The counterbalancing table is the real one from METHODOLOGY_NOTES.md and
%     analysis/design/design_table.csv, NOT the earlier placeholder, which was itself
%     confounded (it held pool constant with condition -- the exact pilot defect).
%   * The OPTT, the SAGAT display-blanking freezes, and the tablet map reconstruction
%     task are OUT. They are absent from the administered forms, so they are absent here.
%   * The post-trial section 3 battery is NOT the SART and is never called that; see
%     score_trial_sa() in analysis/R/10_surveys.R.
% =====================================================================================

\section{User Study Methodology}

\subsection{Participants and Ethical Context}

A total of 24 participants were recruited from a university campus and
surrounding community for a within-subjects, repeated-measures experiment.
Sample size was constrained by the field-deployment logistics of the wide-area
outdoor apparatus: each session required a supervised two-part evening slot
within a narrow dusk-to-night lighting window, at a fixed physical site. We
therefore report a sensitivity analysis in place of an a priori power
calculation. With $N=24$ and the within-subjects contrast of primary interest
(wireframe ON vs.\ OFF), G*Power 3.1 indicates 80\% power ($\alpha=.05$,
two-tailed) to detect effects of $d_z \geq 0.60$, corresponding to $f \approx
0.30$ in the repeated-measures ANOVA parameterization. This study is thus
powered for medium-to-large effects of the kind reported in analogous
wide-area AR spatial cognition work \cite{kim_investigating_2022,
kumaran_impact_2023}; smaller effects may go undetected, and we treat null
results accordingly in Section~\ref{sec:discussion}. Because each participant
contributes four trials (two per wireframe state), trial-level mixed-effects
models over the resulting 96 trials retain more precision than the
participant-mean contrast this sensitivity analysis assumes.

Participant screening and physical safety procedures were strictly governed by
an approved Institutional Review Board (IRB) protocol. Screening covered
uncorrected sight conditions, history of motion sickness or balance disorders,
photosensitive epilepsy, and pregnancy. To control for baseline spatial
ability, participants completed the Santa Barbara Sense of Direction Scale
(SBSOD) prior to trials.

\subsection{Task Design and Counterbalancing}

Participants navigated the physical courtyard to locate 20 virtual targets
embedded among distractor objects. All 20 targets were identical, present from
the start of the trial, and collectable in any order; the task was therefore
free search rather than a prescribed route. Target collection required a
4.0-second gaze dwell time. Targets spawned only within authored walkable
areas. Each of the four object layout pools carries a fixed generation seed, so
every participant assigned to a given pool encounters an identical layout, and
the recall objects for that pool are spawned alongside the targets. Those
recall objects are the distractors: they are everyday scene objects stratified
by physical size, and they serve double duty as visual clutter during the
trial and as the studied items of the recognition test that follows it
(Section~\ref{sec:recall}).

Participants also wore a printed ArUco marker on the wrist, on which the system
renders a wrist-mounted navigation map of the site. The map is a deliberate
element of the apparatus rather than an incidental display: consulting it is
time spent not looking at the environment, so it is measured separately
throughout (Section~\ref{sec:metrics}).

The independent variable (wireframe ON vs.\ OFF) was evaluated using a
rigorously counterbalanced within-subjects design. A previous pilot perfectly
confounded object layout pools with the wireframe condition, rendering initial
results uninterpretable. To fix this, we structured a four-group Latin square
matrix utilizing four distinct object layouts (Pools 1--4), shown in
Table~\ref{tab:counterbalance}. Participants were allocated equally across the
four groups ($n=6$ per group), and group assignment fully determines every
trial's pool and wireframe state. The design satisfies four balance
properties: each pool appears twice with the wireframe and twice without; each
group encounters each pool exactly once; each pool appears exactly once in
each trial position; and each trial position is two ON and two OFF across
groups. Pool is therefore orthogonal to both condition and order.

Daylight cannot be randomized; it falls monotonically through the evening
sessions. Therefore, the wireframe state strictly alternated within a group
(e.g., ON/OFF/ON/OFF), evenly sampling both wireframe states across the
lighting gradient. Groups 1 and 3 run the opposite phase to Groups 2 and 4, so
the small residual association between condition and elapsed session time
cancels across the sample. Lighting is reported strictly as a measured
covariate, never as a manipulation.

\begin{table}[t]
  \caption{Four-group Latin square counterbalancing matrix. Each cell gives the
    object pool and wireframe state for that trial. Trials 1--2 fall in the
    Dusk block and trials 3--4 in the Night block.}
  \label{tab:counterbalance}
  \begin{tabular}{lllll}
    \toprule
    Group & Trial 1 & Trial 2 & Trial 3 & Trial 4\\
    \midrule
    G1 & P1 / OFF & P2 / ON  & P3 / OFF & P4 / ON\\
    G2 & P2 / ON  & P1 / OFF & P4 / ON  & P3 / OFF\\
    G3 & P4 / OFF & P3 / ON  & P2 / OFF & P1 / ON\\
    G4 & P3 / ON  & P4 / OFF & P1 / ON  & P2 / OFF\\
    \bottomrule
  \end{tabular}
\end{table}

\subsection{Object Recall Task}
\label{sec:recall}

Each pool carries two fixed lists: 15 recall objects, which are spawned
alongside that pool's targets during the trial, and 15 lures, which are never
spawned in that pool. Immediately after each trial, participants completed a
30-item confidence-graded recognition test comprising exactly those two lists,
the 15 objects present in that trial and the 15 absent from it. For each item,
participants selected one of three responses: confident they saw the object,
unsure whether they saw it, or confident they did not see it. The answer key
is therefore defined per pool and item position, yielding 120 keyed rows
across the four pools.

\paragraph{Size classes.} Every object in the recall inventory belongs to one
of three size classes: small, medium, or large. Classes were assigned by
direct comparison rather than by a numeric threshold. All 128 candidate props
were instantiated together in a single staging scene at the scale at which
they render during a trial, and arranged there in three parallel rows, one per
class, so that each object was judged against the full set of alternatives at
true size and the resulting assignment is preserved as a property of the scene
rather than as an undocumented decision. A threshold on a single dimension was
rejected because apparent size in the headset follows an object's whole
silhouette rather than any one axis. A slender object such as a light stand or
a structural beam is long but presents very little of itself at the distances
this site affords, and a rule keyed to its longest dimension would file it
alongside vending machines. Comparison in the staging scene resolves those
cases as the participant sees them, which is also why two variants of the same
kind of asset can fall in different classes.

\paragraph{Matching the lures on size.} Each pool's recall list holds five
objects of each class, and each pool's lure list holds five of each class, so
within every pool the present and absent items are matched on size. The four
recall lists and four lure lists together draw 120 objects from the staging
scene, 40 in each class, evenly split between present and absent; the
remaining candidates were held as unused spares.

That matching is a deliberate control rather than a by-product of sampling.
Size governs whether an object is resolvable at the distances this site
affords, so an inventory in which the present objects were systematically the
larger ones would let a participant discriminate seen from unseen on size
alone. Hits and false alarms would then follow a property of the stimulus set
rather than memory for the environment, and $d'$ would overstate recognition
in both conditions. Stratifying both lists closes that route to a correct
answer. It also makes size an analysable factor rather than a nuisance, since
hits and false alarms can be tabulated per class, and any effect of the
wireframe on encoding can be examined separately for objects that are easy and
hard to see from a distance.

Both lists are therefore authored under this constraint rather than sampled at
random. The pool seed fixes \emph{where} a pool's objects are placed
(Section~\ref{sec:layouts}); it does not choose \emph{which} objects the pool
contains, and the size composition of both lists is identical across the four
pools.

Two scoring rules were preregistered and both are computed. Under the
\emph{strict} rule, which is primary, only an unambiguous ``confident saw''
counts as an affirmative response; under the \emph{lenient} rule, ``unsure''
is folded into the affirmative. Reporting both demonstrates that the
conclusion does not hinge on how uncertainty is treated. Recognition
performance is summarised using signal detection theory---$d'$, criterion $c$,
and $A'$---with the loglinear correction \cite{hautus_corrections_1995}
applied to every trial rather than only to degenerate cells. With 30 items a
hit rate of 1.00 is common and the uncorrected $d'$ would be infinite;
correcting only the broken cells would bias exactly those cells. An item-level
generalised linear mixed-effects model complements the aggregate SDT measures.

\paragraph{Incidental encoding on Trial 1.} Participants were not informed of
the recall task before their first trial. Encoding on Trial 1 is therefore
\emph{incidental}, whereas encoding on Trials 2--4 is \emph{intentional}, as
participants then know a recognition test follows. To capture this asymmetry,
participants additionally rated their confidence in their first recall test on
a 7-point scale. This distinction is a property of trial position rather than
of condition, and the counterbalancing matrix
(Table~\ref{tab:counterbalance}) places two wireframe-ON and two wireframe-OFF
groups in trial position 1, so encoding intentionality is orthogonal to the
independent variable. We nevertheless model trial position explicitly and
report Trial 1 recall separately, since collapsing incidental and intentional
encoding into a single recall score would inflate within-participant variance.

\subsection{Measures}

\textbf{Pre-study (once, before trials).} Demographics (age band, gender),
self-rated AR/VR comfort (5-point), prior AR/VR experience, cumulative usage
band, and the safety screening described above. Baseline spatial ability was
measured with the 15-item SBSOD (7-point), of which 8 items are reverse-coded.

\textbf{Post-trial (after each of the four trials).} Raw NASA-TLX (six items,
0--100 in increments of 5; the Performance item is \emph{not} reverse-coded,
as it is already worded higher-is-worse); a 10-item trial-level situation
awareness battery (7-point) combining a subset of SART dimensions
(complexity, spare mental capacity, concentration, division of attention,
familiarity) with items probing perception, comprehension, and projection of
the environment; a four-item memory and brightness block covering self-rated
recall success, perceived change in recall performance over time, and two
bipolar brightness items scored as deviation from the scale midpoint; a
purpose-designed 7-point item assessing perceived visual clutter; and the
five-item MEC-SPQ Spatial Situational Model subscale (5-point).

\textbf{Post-study (once, after all trials).} The 10-item SART in Taylor's
three-factor form (Demand, Supply, Understanding), scored as $U - (D - S)$;
the 16-symptom Simulator Sickness Questionnaire (SSQ); a five-item wireframe
helpfulness block (7-point); a repeat of the memory and brightness block at
study level; a four-item spatial awareness block covering disorientation
during and after headset use and perceived collision risk; and open-ended
items on perceived study purpose, noticed layout patterns, and how the
wireframe did or did not help.

Because the post-study measures are collected once per participant, they
cannot enter the condition model and are reported as descriptive and
correlational context only. The trial-level situation awareness battery is
consequently the only situation-awareness measure that varies with wireframe
state, and it is reported both as its three subscale means and at the level of
individual items. We report it separately from the post-study SART throughout;
the two are distinct instruments on different scales and are never pooled or
compared.

\subsection{Experimental Procedure}

To mitigate the CVIP memory exhaustion crash identified in hardware testing,
the study was physically split into two parts separated by a headset reboot.

\begin{enumerate}
  \item \textbf{Onboarding:} Informed consent, safety screening, and the
    pre-study survey (demographics and SBSOD). Participants were briefed on
    the target search task only; the object recall task was deliberately not
    disclosed, preserving incidental encoding for Trial 1.
  \item \textbf{Part 1 (Dusk):} Trials 1 and 2. After each trial,
    participants completed the 30-item recall test followed by the post-trial
    battery (NASA-TLX, situation awareness, memory and brightness, visual
    clutter, MEC-SPQ). The Trial 1 recall test additionally included the
    7-point confidence rating described above.
  \item \textbf{Inter-part break:} Participants relocated indoors while the
    headset was rebooted and recharged.
  \item \textbf{Part 2 (Night):} Following the reboot, markers were
    re-scanned and participants completed Trials 3 and 4, each again followed
    by the recall test and the post-trial battery.
  \item \textbf{Post-study and debrief:} The post-study survey (SART, SSQ,
    wireframe helpfulness, memory and brightness, spatial awareness, and
    open-ended items), followed by debrief.
\end{enumerate}

% =====================================================================================
% Analytical Strategy -- corrected against the implemented pipeline.
%
% REMOVED, because the pipeline does not produce them and no task collects them:
%   * egocentric pointing angular error
%   * relative object location recall distance
%   * feature-binding accuracy
%   * "fixations within 2 degrees of rendered edges" (AOI attribution is collider-based)
%   * OPTT as a covariate
%   * pupil diameter, in every form -- as a workload measure, as a caveated secondary,
%     and as a stated hardware limitation. The section does not mention it at all, by
%     decision. It is not measured: the ML2 OpenXR path returns one combined gaze pose
%     and no pupillometry, so EyeAndHeadTracker writes pupilDiameterMm = -1 (see the
%     comment at ~line 967) and the reduction maps that sentinel to NA. Do not reinstate
%     the topic because a reviewer asks for pupillometry; the reasoning is kept in
%     METHODOLOGY_NOTES.md and analysis/README.md, not in the paper.
%
% RESTORED, now that it is actually implemented: the trajectory / ideal-path efficiency
% ratio. GemOptimalPathCalculator.cs writes analysis/design/ideal_paths.csv (exact
% Held-Karp open tour per pool) and 12_join.R divides head_path_xz_m by it. Both terms
% are planar so they are measured in the same plane.
%
% The denominator is now WALKABLE, not straight-line: distances come from Dijkstra over a
% 0.5 m grid ray-sampled from the site geometry (DistanceMode.WalkableGrid), so the tour no
% longer passes through buildings. The old "straight-line, therefore a lower bound" caveat
% is GONE from the text because it is no longer true -- do not reinstate it. Three limits
% replace it and must stay: the 0.5 m / 8-connected grid approximation, the hand-authored
% gap-bridging planes, and known target positions. If the mode is ever switched back to
% StraightLine, this paragraph has to change with it; the CSV's `method` column records
% which mode produced each row, so the two can always be checked against each other.
%
% The walkable region and the spawn region are DELIBERATELY separate objects in the scene:
% RandomSpawner.targetPlanes is frozen (its area and list length feed the layout seeds --
% changing it regenerates every pool), and walk-only surfaces live on a WalkableAreaSet
% component. Do not merge them.
%
% UNRESOLVED -- decide before any real reduction run. The gaze screening threshold is
% stated as 50% valid in METHODOLOGY_NOTES.md but implemented as "exclude above 20%
% missing" in 00_config.R (max_missing_gaze = 0.20). They are not compatible: 20%
% missing requires 80% valid, and the outdoor pilots returned 62.1% and 70.4% valid --
% so as configured, both pilot sessions would be excluded from every gaze model. The
% text below uses the PREREGISTERED 20% because that is what the code does. If the 50%
% floor is what is intended, change cfg$max_missing_gaze to 0.50 and this paragraph
% together; do not change the prose alone.
% =====================================================================================

\section{Analytical Strategy and Evaluation Metrics}

To isolate scaffolding effects from spatial-ability variance and
object-property confounds, we structured a mixed-methods statistical pipeline
using the R programming language (\texttt{lme4}, \texttt{lmerTest}, and
\texttt{emmeans} packages). Every threshold below is fixed in a single
configuration file, written into the run log, and applied automatically; the
reduction pipeline is validated end-to-end on simulated data before real data
is touched.

\subsection{Telemetry and Gaze Processing}

Eye and head telemetry was logged at approximately 60\,Hz via a non-blocking
dual-stream logging architecture. To ensure data survival against crashes,
look events were streamed append-only to their own \texttt{.ndjson} file,
while the JSON summary is rewritten in full on each autosave. Durations were
summed from inter-row intervals rather than counted in rows, so that dropped
frames are not silently rescaled: a 59-row second is not $59/60$ of a second.
Inter-row intervals exceeding 0.2\,s---twelve dropped frames at 60\,Hz---are
treated as unobserved time, excluded from every duration and breaking any
dwell or eye event that spans them. This matters because no rows are written
while a trial is operator-paused, so a pause appears as one enormous interval;
counting it would credit the participant with minutes of dwell on whatever
they happened to be looking at when the trial stopped. Quantities that were
not measured are written as empty cells rather than zeros, so that, for
example, saccade peak velocity during a fixation cannot drag a study-wide mean
toward zero.

Gaze is attributed to an area of interest only when eye tracking is active,
the gaze sample is valid, and confidence meets the floor recorded in that
trial's own file header. Looks at the same AOI separated by less than 0.10\,s
are merged, so that a single blink mid-look is not scored as two looks, and
looks shorter than 0.10\,s are discarded as ray-crossings. These three
values---the exclusion threshold, the merge window, and the minimum look
duration---were preregistered before any real data was examined.

Trials in which more than 20\% of trial time lacked usable gaze were excluded
from eye-tracking models but expressly retained for survey and recall
analyses, since eye-tracker quality is uninformative about whether a
participant completed a questionnaire. Exclusions are reported as an
automatically generated CONSORT-style flow. Ambient infrared interference
outdoors is the limiting factor on tracking quality: indoor sessions run
around 89\% valid gaze, whereas outdoor pilot sessions returned 62.1\% and
70.4\%. We report that gap as a finding in its own right, as it constrains
what wide-area outdoor eye-tracking studies on this hardware can measure.

\subsection{Evaluation Metrics}
\label{sec:metrics}

Measures are grouped by the level at which they vary, because that determines
which may enter a condition model at all.

\begin{itemize}
  \item \textbf{Task performance (per trial):} target acquisition counts and
    inter-object time deltas.
  \item \textbf{Locomotion and scanning (per trial):} ground-plane head path
    length, cumulative yaw travel, and a search efficiency ratio---distance
    actually walked divided by the length of the shortest route that would
    have visited all 20 targets. Because the targets are identical and
    collectable in any order, that ideal route is an open travelling-salesman
    tour over the pool's 20 generated positions (Section~\ref{sec:layouts}),
    solved exactly by Held--Karp at authoring time and stored per pool.

    The tour is measured across the walkable surface rather than in straight
    lines, so it does not pass through buildings. Inter-target distances are
    shortest paths over a 0.5\,m grid sampled from the site geometry by
    downward ray cast---the same test that constrains where objects may
    spawn---giving 22\,117 traversable cells covering
    5529\,m\textsuperscript{2}. Shortest paths between cells are computed by
    Dijkstra's algorithm on an 8-connected lattice, and the resulting distance
    matrix is what Held--Karp then optimises over. The traversable region is
    the union of the spawn surface with a small number of hand-placed planes
    bridging gaps left by the scan: places a participant can walk but where no
    object is ever placed. Those bridging surfaces are authored rather than
    measured, which we note as a limitation; they shorten the ideal tour by
    8\% in Pool 1 and 6\% in Pool 3 and leave Pools 2 and 4 unchanged.
    Constraining the route to walkable ground lengthens the tour by 10--13\%
    relative to a straight-line tour over the same visiting order
    (Table~\ref{tab:pools}), and it reorders the pools by tour length, so the
    constraint is not a per-pool rescaling of the straight-line measure.

    Three properties govern how the ratio is read. The grid discretises the
    surface at 0.5\,m and connects cells in eight directions, so a diagonal
    run is measured up to roughly 8\% long; this bias is identical for every
    trial in a pool and therefore cancels in the between-condition comparison
    the ratio is used for. The denominator assumes the target positions are
    known in advance, which the participant's are not, since finding them is
    the task---so the ratio is always well above unity and is interpreted as a
    relative index compared across conditions within a pool, never as an
    absolute claim about how far a participant overshot. Finally, path length
    is measured in the ground plane so that numerator and denominator share a
    plane, and because the vertical component of a head-mounted sensor's
    trajectory accumulates gait rather than route choice; the same tours
    measured in three dimensions are within 0.5\% of their planar lengths,
    confirming that the routes are effectively flat.
  \item \textbf{Gaze dynamics (per trial):} the primary visual-attention
    outcome is the share of usable gaze time resting on non-target scene
    geometry, time-weighted from per-sample intervals rather than frame
    counts, so that the frame-rate cost of rendering the wireframe cannot move
    it by itself. The secondary outcome is the number of distinct pieces of
    scenery fixated per minute of tracked time, since a raw distinct-object
    count is uninterpretable when trials differ in length. Reported alongside:
    time attributed to no AOI, number of switches between AOIs, and the
    per-AOI dwell breakdown.
  \item \textbf{Map consultation (per trial):} dwell on the wrist-mounted
    navigation map is scored under its own category and is deliberately
    \emph{excluded} from the scenery outcome above rather than pooled into it.
    The map is a UI panel the participant chooses to consult, not scenery they
    noticed, and it is a single object attracting very long dwells: in pilot
    data it accounted for 33.9\,s of the 35.3\,s of non-target dwell in one
    trial and 46.5\,s of 56.0\,s in the next. Left in, the primary outcome
    would read as ``time spent consulting the map,'' a different construct
    that moves in the opposite direction, since looking at the map is time
    spent not looking at the world. It is reported separately as a covariate
    on how much of the environment a participant examined.
  \item \textbf{Recognition memory (per trial):} $d'$, criterion $c$, and
    $A'$ from the 30-item recall test, under both the strict and lenient
    response policies.
  \item \textbf{Subjective, per trial:} NASA-TLX, the trial-level situation
    awareness battery, perceived visual clutter, MEC-SPQ, and the trial-level
    memory and brightness items.
  \item \textbf{Subjective, per participant:} SART, SSQ total severity,
    wireframe helpfulness, and the study-level spatial awareness items. These
    are collected once and therefore \emph{cannot} enter the condition model;
    they are reported descriptively and correlationally.
\end{itemize}

\subsection{Statistical Modeling}

Recall accuracy was scored using signal detection theory ($d'$, criterion,
$A'$) applying a loglinear correction across all trials rather than only to
degenerate cells. Item-level recognition responses were modeled using binomial
generalized linear mixed models (GLMMs) with crossed random effects for
participant and item. Continuous per-trial metrics were analyzed using linear
mixed models (LMMs). For visual attention (H4), an LMM regressed the
time-weighted proportion of usable gaze resting on non-target scene geometry
against wireframe state, with map dwell excluded from that proportion and
entered separately, and with distinct scenery objects fixated per minute as
the secondary outcome. The per-AOI breakdown separates attention to
scaffolding geometry from attention to task objects. Search efficiency was
modeled on the same terms, using the ratio of walked to ideal path length;
note that since the denominator is constant within a pool, the ratio is a
per-pool rescaling of path length and is reported for comparability rather
than as an independent test. Subjective perceived clutter ($H_{D3}$) was
evaluated using ordinal logistic regression on the 7-point visual clutter
item, with the NASA-TLX Mental Demand item acting as a secondary convergent
measure.

Across all models, wireframe state was the fixed effect of interest; SBSOD
score was included as a covariate to control for baseline spatial ability;
object pool and trial position were included as design factors; and
participant was a random intercept. Trial position additionally carries the
incidental-versus-intentional encoding distinction described in
Section~\ref{sec:recall}, which is why it is modeled explicitly rather than
absorbed. Significance was assessed via Type-II likelihood-ratio tests with
Bonferroni-corrected post-hoc comparisons.
